\section{Conclusion}
\label{sec:conclusion}

We presented \methodname{}, a bilevel framework in which an outer loop
autonomously discovers and injects structural changes into an inner autoresearch
loop.
The key result is that Level~2---an autonomous mechanism research agent that
writes and loads Python code at runtime---produces a 5$\times$ improvement in
\valbpb{} over the Level-1 baseline on Karpathy's GPT pretraining benchmark,
using the same DeepSeek model at every level.
Level~1.5 alone does not reliably improve over the pure inner loop, confirming
that parameter-level adjustment is insufficient to escape the LLM's
deterministic search patterns; structural modification of the loop is required.
Group~D (Level~1 + Level~2, without Level~1.5) achieves $-0.034 \pm 0.031$,
confirming that Level~2 is the key driver of improvement and that Level~1.5
is not essential when mechanism research is present across four experimental groups.

Level~2's mechanisms---Tabu Search Manager, Multi-Scale Bandit Proposer, and
Systematic Orthogonal Exploration---were drawn from three distinct algorithmic
domains and discovered without human specification, across three independent
repeats (a fourth domain, Bayesian optimization, was attempted but reverted
due to a missing dependency).
These mechanisms succeeded by forcing exploration of a search direction
(\texttt{TOTAL\_BATCH\_SIZE} reduction) that the LLM's default priors
systematically avoid.

Code, experiment logs, and all Level~2 generated mechanisms are available at
\url{https://github.com/EdwardOptimization/Bilevel-Autoresearch}.
The paper is available at \url{https://aixiv.science/abs/aixiv.260323.000006}.

The core principle is validated: \emph{autoresearch can research itself.}
The outer loop need not be human-designed; the same model that runs the inner
loop can generate it, discovering structural improvements that previously
required a human researcher to write.
